% Chapter 4: Experiments — datasets, leakage audit, separability, CascadeLP/SVM results
\selectlanguage{greek}


\chapter{Πειράματα}
\label{ch:experiments}

Το παρόν κεφάλαιο αξιολογεί το πλαίσιο \en{CascadeLP} του
Κεφαλαίου~\ref{ch:methodology} σε δύο σύνολα υπερσυνδέσμων της \en{Wikipedia}.
Το \en{DSAA 2023} επιτρέπει άμεση σύγκριση με τις λύσεις του διαγωνισμού, αλλά ο
έλεγχος ακεραιότητας της \S\ref{sec:hub-artifact} περιορίζει ουσιωδώς την
ερμηνεία των σχεδόν τέλειων βαθμολογιών του. Για τον λόγο αυτό κατασκευάζεται
δεύτερο σύνολο από επίσημα αρχεία της Απλής Αγγλικής \en{Wikipedia}, στο οποίο
το ίδιο αρχιτεκτονικό σχήμα επανεκπαιδεύεται και αξιολογείται με αυστηρότερο
πρωτόκολλο. Τα δύο πειράματα απαντούν σε συμπληρωματικά ερωτήματα: το πρώτο
μετρά τη συμπεριφορά στο καθιερωμένο πεδίο αξιολόγησης, ενώ το δεύτερο παρέχει
την κύρια ένδειξη για τη χρησιμότητα της αρχιτεκτονικής.


\section{Σύνολα Δεδομένων}
\label{sec:datasets}

\subsection{\en{Wikipedia DSAA 2023}}

Το πρώτο σύνολο είναι το \en{Wikipedia DSAA 2023}
(\S\ref{sec:dsaa-dataset}), με \TrainPairs{} επισημασμένα ζεύγη εκπαίδευσης,
\TestPairs{} ζεύγη ελέγχου και \GraphNodesTotal{} άρθρα. Διατηρείται επειδή
αποτελεί το αρχικό πεδίο ανάπτυξης του \en{CascadeLP} και διαθέτει δημόσιο
\en{leaderboard}. Οι βαθμολογίες του αναφέρονται ως αποτελέσματα στο συγκεκριμένο
σύνολο, χωρίς να θεωρούνται από μόνες τους απόδειξη εκμάθησης υπερσυνδέσμων.

\subsection{Ανεξάρτητο Σύνολο Απλής Αγγλικής \en{Wikipedia}}
\label{subsec:wiki-fresh-construction}

Το δεύτερο σύνολο κατασκευάζεται από τους επίσημους πίνακες \texttt{\en{page}},
\texttt{\en{linktarget}} και \texttt{\en{pagelinks}} της Απλής Αγγλικής
\en{Wikipedia}. Από το πλήρες γράφημα συλλέγεται μέσω \en{BFS}, με αφετηρία το
άρθρο \en{«Computer science»}, συνεκτικό υπογράφημα \WfCrawlNodes{} κόμβων και
\WfCrawlEdges{} πραγματικών μη κατευθυνόμενων ακμών. Οι θετικές εγγραφές είναι
οι πραγματικοί υπερσύνδεσμοι του υπογραφήματος. Για κάθε θετική ακμή
δειγματοληπτείται ένα διαφορετικό, ομοιόμορφα τυχαίο μη συνδεδεμένο ζεύγος από
τον ίδιο πληθυσμό κόμβων.

Το κείμενο ανακτάται από το σύνολο \en{wikimedia/wikipedia} με διαμόρφωση
\texttt{\WfTextConfig{}}. Η ακριβής ταύτιση τίτλου βρίσκει κείμενο για
\WfTextMatched{} από τους \WfTextWanted{} κόμβους (\WfTextMatchPct\%), ενώ
\WfTextEmptyN{} κόμβοι διατηρούν κενό κείμενο. Η απουσία αυτή αποτελεί
περιορισμό των κειμενικών συγκρίσεων και δεν θεωρείται εκ προοιμίου αμελητέα.
Η κατασκευή αποφεύγει το συγκεκριμένο μοτίβο συγκέντρωσης αρνητικών του
\en{DSAA 2023}, αλλά η επιλογή ενός σπόρου \en{BFS} και εύκολων τυχαίων
αρνητικών επηρεάζει τη δυσκολία του νέου συνόλου και συζητείται στους
περιορισμούς (\S\ref{sec:limitations-cost}).


\section{Διάταξη Πειραμάτων}
\label{sec:experimental-setup}

\subsection{Διαιρέσεις Δεδομένων}
\label{subsec:data-splits}

Στο \en{DSAA 2023}, το \texttt{\en{train.csv}} περιέχει \TrainPairs{} ζεύγη.
Εξαιρούμε τα \SelfLoopsTrain{} ζεύγη
αυτοβρόχων --- τα οποία επιλύονται ξεχωριστά στο Επίπεδο~0 (\S\ref{sec:cascade-routing})
--- και διαιρούμε τα υπόλοιπα \TrainPairsNoSelf{} ζεύγη σε σύνολο εκπαίδευσης και σύνολο
επικύρωσης με αναλογία 80\%/20\%, στρωματοποιημένη ως προς την ετικέτα, με σταθερό
σπόρο τυχαιότητας (\texttt{\en{random\_state=42}}). Αυτό δίνει \TrainSplitSize{} ζεύγη
εκπαίδευσης και \ValSplitSize{} ζεύγη επικύρωσης, τα ίδια μεγέθη σε όλα τα μοντέλα που
αναφέρονται σε αυτό το κεφάλαιο, ώστε οι συγκρίσεις μεταξύ τους να είναι ισοδύναμες.
Ο διαχωρισμός γίνεται μία φορά και επαναχρησιμοποιείται· κανένας ταξινομητής δεν
βλέπει τα ζεύγη επικύρωσης κατά την εκπαίδευσή του.

Το αρχικό πείραμα \en{DSAA} έχει έναν σημαντικό περιορισμό: το δομικό γράφημα
κατασκευάζεται από το σύνολο των θετικών ακμών ολόκληρου του
\texttt{\en{train.csv}}, πριν από τον διαχωρισμό εκπαίδευσης/επικύρωσης, όπως και τα
τέσσερα ευρετικά χαρακτηριστικά που
υπολογίζονται πάνω του (Κεφάλαιο~\ref{ch:methodology}). Έτσι, οι παρακρατημένες
θετικές ακμές έχουν ήδη επηρεάσει το γράφημα, τους βαθμούς και ενδεχομένως τα
δομικά χαρακτηριστικά τους. Το σύνολο είναι παρακρατημένο από την εκπαίδευση των
ταξινομητών, αλλά δεν είναι ανεξάρτητο ως προς την κατασκευή του γραφήματος. Για τον λόγο αυτό οι τιμές επικύρωσης του
\en{DSAA} αναφέρονται ως αποτελέσματα του αρχικού πρωτοκόλλου και δεν
χρησιμοποιούνται ως η κύρια απόδειξη εγκυρότητας του \en{CascadeLP}.
Ο \texttt{\en{SvmClassifier}}
(\S\ref{sec:svm-baseline}) εκπαιδεύεται σε ένα επιπρόσθετο, στρωματοποιημένο
δείγμα \SvmSampleSize{} ζευγών από το σύνολο εκπαίδευσης, λόγω του υπολογιστικού κόστους
$O(n^2\text{--}n^3)$ του \en{SVC} πυρήνα ακτινικής βάσης.

Το \texttt{\en{test.csv}} (\TestPairs{} ζεύγη) παραμένει εντελώς ανεξάρτητο από αυτή τη
διαδικασία και δεν διαθέτει δημοσιευμένες πραγματικές ετικέτες τοπικά· η αξιολόγηση
σε αυτό γίνεται αποκλειστικά μέσω υποβολής στον πίνακα κατάταξης \en{Kaggle} του
διαγωνισμού (\S\ref{sec:cascade-results}).

\begin{figure}[h!]
    \centering
    \tikzset{
        box/.style={rectangle, draw, rounded corners, minimum width=3.4cm,
            minimum height=1cm, align=center},
        note/.style={rectangle, draw, dashed, rounded corners, minimum width=3.6cm,
            minimum height=1.1cm, align=center, fill=black!5, font=\small},
        flow/.style={-{Stealth}, thick},
    }
    \resizebox{\linewidth}{!}{%
    \begin{tikzpicture}[node distance=1.0cm and 1.6cm]
        \node[box] (raw) {\texttt{\en{train.csv}}\\\TrainPairs{} ζεύγη};
        \node[box, below=of raw] (noloop) {Εξαίρεση αυτοβρόχων\\($-$\SelfLoopsTrain{})};
        \node[box, below=of noloop] (split938) {\TrainPairsNoSelf{} ζεύγη};
        \node[box, below=1.6cm of split938, xshift=-2.4cm] (tr) {Σύνολο Εκπαίδευσης\\\TrainSplitSize{} (80\%)};
        \node[box, below=1.6cm of split938, xshift=2.4cm] (val) {Σύνολο Επικύρωσης\\\ValSplitSize{} (20\%)};

        \node[note, right=4.6cm of raw] (graph) {Δομικό γράφημα \&\
            (όλα τα \TrainPairs{}, \textit{πριν} τη διαίρεση)};
        \node[box, below=1.6cm of graph] (test) {\texttt{\en{test.csv}}\\\TestPairs{} ζεύγη\\(ανεξάρτητο)};
        \node[note, below=1.0cm of test] (kaggle) {Αξιολόγηση μόνο μέσω\\πίνακα κατάταξης \en{Kaggle}};

        \draw[flow] (raw) -- (noloop);
        \draw[flow] (noloop) -- (split938);
        \draw[flow] (split938.south) -- ++(0,-0.5) -| node[pos=0.25, above]{στρωματοπ.} (tr.north);
        \draw[flow] (split938.south) -- ++(0,-0.5) -| node[pos=0.25, above]{\en{seed=}\Seed{}} (val.north);
        \draw[flow, dashed] (raw.east) -- (graph.west);
        \draw[flow] (test) -- (kaggle);
    \end{tikzpicture}%
    }
    \caption{Διαίρεση δεδομένων: το δομικό γράφημα υπολογίζεται στο πλήρες \texttt{\en{train.csv}} πριν από τον διαχωρισμό
    (διακεκομμένο βέλος), ενώ κάθε ταξινομητής εκπαιδεύεται αποκλειστικά στο 80\%
    και αξιολογείται στο παρακρατημένο 20\%. Το \texttt{\en{test.csv}} δεν συμμετέχει σε
    καμία φάση εκπαίδευσης.}
    \label{fig:data-split}
\end{figure}

Στο ανεξάρτητο σύνολο χρησιμοποιείται ομαδοποιημένη διαίρεση, ώστε το ίδιο μη
κατευθυνόμενο ζεύγος να μην μπορεί να εμφανιστεί και στις δύο πλευρές. Το δομικό
γράφημα κατασκευάζεται αποκλειστικά από τις θετικές ακμές του τμήματος
εκπαίδευσης. Κατά τον υπολογισμό χαρακτηριστικών μιας θετικής ακμής-στόχου, η
ακμή αφαιρείται προσωρινά από το γράφημα, αποτρέποντας την άμεση χρήση της
ετικέτας που επιχειρεί να προβλέψει το μοντέλο. Η διαίρεση δίνει
\WfTrainSplitSize{} ζεύγη εκπαίδευσης και \WfValSplitSize{} ζεύγη επικύρωσης.

\subsection{Μετρικές Αξιολόγησης}
\label{subsec:evaluation-metrics}

Πρωτεύουσα μετρική είναι το \en{Macro F1-score}: ο μέσος όρος του \en{F1-score} ανά
κλάση, υπολογισμένος χωρίς στάθμιση ως προς το πλήθος δειγμάτων κάθε κλάσης.
Η
επιλογή αυτή δεν είναι σχηματική --- το Κεφάλαιο~\ref{ch:analysis} δείχνει
συγκεκριμένα υποσύνολα (π.χ.\ το εμπιστευτικό υποσύνολο του Επιπέδου~1,
\S\ref{sec:cascade-results}) όπου η απλή ακρίβεια παραμένει άνω του 0{,}99 ενώ το
\en{Macro F1} καταρρέει σε τιμή κοντά στο 0{,}50, ακριβώς λόγω ανισορροπίας κλάσεων
που η ακρίβεια αδυνατεί να συλλάβει.
Ως δευτερεύουσες μετρικές αναφέρουμε το
\en{AUC-ROC}, το \en{Macro F1} περιορισμένο στο διαγνωστικό υποσύνολο μηδενικού
\en{CN} (η ιστορικά ονομασμένη \texttt{\en{cold\_start\_mask}} --- μηδενικοί
κοινοί γείτονες μεταξύ $u$ και $v$) και τον χρόνο συμπερασμού ανά ζεύγος. Η
μάσκα αυτή είναι ευρύτερη από το λειτουργικό \en{cold-start} της δρομολόγησης,
το οποίο απαιτεί απουσία τουλάχιστον ενός άκρου από το γράφημα.
Οι
υπολογισμοί αυτών των μετρικών (\texttt{\en{src\slash{}utils\slash{}metrics.py}}) είναι κοινοί σε
όλα τα μοντέλα αυτού του κεφαλαίου, ώστε οι συγκρίσεις να είναι άμεσα συγκρίσιμες.
Ο \textbf{Αλγόριθμος~\ref{alg:eval-protocol}} περιγράφει το κοινό πρωτόκολλο
αξιολόγησης που εφαρμόζεται σε κάθε μοντέλο αναφοράς και στο \en{CascadeLP}.

\begin{algorithm}[h!]
\caption{Πρωτόκολλο Αξιολόγησης Μοντέλου}
\label{alg:eval-protocol}
\begin{algorithmic}[1]
\Require Ζεύγη εκπαίδευσης $D$ με ετικέτες $y$· οικογένεια χαρακτηριστικών $\varphi$·
    μοντέλο $M$
\Ensure Μετρικές αξιολόγησης στο σύνολο επικύρωσης
\State $D \gets D \setminus \{(u,v) \in D : u = v\}$ \Comment{εξαίρεση αυτοβρόχων}
\State $(D_{\text{tr}}, D_{\text{val}}) \gets \text{StratifiedSplit}(D, y,
    \text{test\_size}=\ValSizeFrac{}, \text{seed}=\Seed{})$
\State $X_{\text{tr}} \gets \varphi(D_{\text{tr}})$;\quad
    $X_{\text{val}} \gets \varphi(D_{\text{val}})$
    \Comment{η κατασκευή του γραφήματος καθορίζεται από το εκάστοτε πρωτόκολλο}
\State $M.\text{fit}(X_{\text{tr}}, y_{\text{tr}})$
\State $s_{\text{val}} \gets M.\text{predict\_proba}(X_{\text{val}})[:,1]$
\State $\hat{y}_{\text{val}} \gets \mathbb{I}[s_{\text{val}} \geq 0{,}5]$
\State $\text{macroF1} \gets \text{F1}(y_{\text{val}}, \hat{y}_{\text{val}}, \text{average}=\text{macro})$
\State $\text{auc} \gets \text{AUC-ROC}(y_{\text{val}}, s_{\text{val}})$
\State $\text{csF1} \gets \text{F1}\big(y_{\text{val}}[\texttt{cold\_start\_mask}],\;
    \hat{y}_{\text{val}}[\texttt{cold\_start\_mask}]\big)$
\If{$M$ είναι \en{CascadeLP}}
    \State $\text{difficulty} \gets \text{label\_difficulty}(D_{\text{val}})$
    \Comment{\S\ref{sec:forensics-protocol}}
    \State $\text{tier} \gets$ επίπεδο που επέλυσε κάθε ζεύγος του $D_{\text{val}}$
    \State υπολογισμός πίνακα διασταύρωσης $(\text{tier} \times \text{difficulty})$
    \Comment{\texttt{tier\_difficulty\_breakdown}}
\EndIf
\State \Return $(\text{macroF1}, \text{auc}, \text{csF1})$
\end{algorithmic}
\end{algorithm}

\subsection{Μοντέλα Αναφοράς}
\label{subsec:baselines}

Κάθε ένα από τα τρία εκπαιδευόμενα επίπεδα του \en{CascadeLP} αξιολογείται επίσης ως
αυτόνομο μοντέλο αναφοράς στην ίδια διαίρεση δεδομένων, ώστε να είναι μετρήσιμη η
συνεισφορά καθενός μεμονωμένα: ο \texttt{\en{StructuralClassifier}} (Λογιστική
Παλινδρόμηση στα τέσσερα ευρετικά χαρακτηριστικά), ο \texttt{\en{TfidfClassifier}}
(Λογιστική Παλινδρόμηση στην ομοιότητα συνημίτονου \en{TF-IDF}), ο
\texttt{\en{PosClassifier}} (\en{Random Forest} στα χαρακτηριστικά μερών του λόγου) και ο
\texttt{\en{EmbeddingClassifier}} (Λογιστική Παλινδρόμηση στην ομοιότητα συνημίτονου
ενσωματώσεων \en{sentence-transformer}).
Προσθέτουμε τον \texttt{\en{SvmClassifier}}
(\en{SVC} πυρήνα ακτινικής βάσης στην ίδια ομοιότητα \en{TF-IDF}, με
$C=\SvmC{}$, $\gamma=$\en{«\SvmGamma{}»}, υποδείγμα \SvmSampleSize{} ζευγών) ως ιστορικό σημείο αναφοράς
κλασικής μηχανικής μάθησης, ακολουθώντας το προηγούμενο των \textcite{alhasan2006link} (\S\ref{sec:svm-baseline}).
Η σύγκριση \en{CascadeLP} έναντι
\en{SVM} αναλύεται στην \S\ref{sec:svm-results}.

\subsection{Λεπτομέρειες Υλοποίησης}
\label{subsec:implementation-details}

Ολόκληρη η αλυσίδα επεξεργασίας υλοποιείται σε \en{Python 3.10} μέσα σε
απομονωμένο περιβάλλον \en{conda} (\texttt{\en{link-\allowbreak{}prediction}}), με \en{PyTorch}~\cite{paszke2019pytorch} σε
διαμόρφωση μόνο \en{CPU}.
Οι ταξινομητές Λογιστικής Παλινδρόμησης, Τυχαίου Δάσους
και \en{SVC} χρησιμοποιούν το \en{scikit-learn}· το γράφημα και οι δομικοί
ευρετικοί δείκτες υπολογίζονται με το \en{NetworkX}~\cite{hagberg2008exploring}· η επισήμανση μερών του λόγου
χρησιμοποιεί το \en{NLTK} (\texttt{\en{punkt}} και
\texttt{\en{averaged\_perceptron\_tagger}}, \S\ref{sec:cascade-routing}); οι
ενσωματώσεις κειμένου παράγονται με το μοντέλο \en{sentence-transformer}
\texttt{\en{\EmbeddingModelName{}}}.
Ο σπόρος τυχαιότητας
\Seed{} διατηρείται σταθερός σε όλες τις διαιρέσεις δεδομένων και σε όλα τα μοντέλα,
ώστε τα πειράματα να είναι αναπαραγώγιμα.
Ο πλήρης πηγαίος κώδικας είναι δημόσια
διαθέσιμος (Κεφάλαιο~\ref{ch:introduction}).


\section{Έλεγχος Διαρροής Δεδομένων}
\label{sec:leakage-results}

Εφαρμόζοντας το πρωτόκολλο ανίχνευσης της \S\ref{sec:forensics-protocol} στα
πλήρη αρχεία \texttt{\en{train.csv}} (\TrainPairs{} ζεύγη) και \texttt{\en{test.csv}}
(\TestPairs{} ζεύγη) του \en{DSAA 2023}, βρίσκουμε:

\begin{itemize}
    \item \textbf{Ακριβής επικάλυψη ζευγών}: \LeakageExact{} ζεύγη (\LeakageExactPct\% του συνόλου ελέγχου).
    \item \textbf{Αντεστραμμένη επικάλυψη ζευγών}: \LeakageReversed{} ζεύγη (\LeakageReversedPct\% του συνόλου
    ελέγχου).
    \item \textbf{Αυτοβρόχοι}: \SelfLoopsTrain{} στο σύνολο εκπαίδευσης (\SelfLoopsTrainPos{} θετικοί, \SelfLoopsTrainNeg{} αρνητικός), \SelfLoopsTest{} στο σύνολο ελέγχου.
    \item \textbf{Ενδο-συνόλου διπλοεγγραφή}: \IntraTrainDuplicates{} ομάδες διπλότυπων μη
    κατευθυνόμενων ζευγών εντός του \texttt{\en{train.csv}} (\IntraTrainDuplicatesRows{} γραμμές), χωρίς
    καμία σύγκρουση ετικέτας μεταξύ διπλότυπων.
\end{itemize}

Τόσο η ακριβής όσο και η αντεστραμμένη επικάλυψη είναι αμελητέες ως ποσοστό του
συνόλου ελέγχου.
Αυτό αποκλείει την απλή επικάλυψη ζευγών ως εξήγηση των σχεδόν τέλειων
βαθμολογιών που αναφέρονται στη βιβλιογραφία (Κεφάλαιο~\ref{ch:background}),
όχι όμως άλλες μορφές διαρροής ή τεχνουργήματα δειγματοληψίας. Η ενδο-συνόλου διπλοεγγραφή, αν και μικρή σε
απόλυτο αριθμό, αποτελεί ένα δευτερεύον κανάλι διαρροής μέσω της τυχαίας
διαίρεσης εκπαίδευσης-επαλήθευσης βάσει γραμμής που χρησιμοποιείται στα
πειράματα του \en{CascadeLP} (\S\ref{sec:cascade-routing}).


\section{Χαρακτηρισμός Διαχωρισιμότητας}
\label{sec:separability-results}

Εφαρμόζουμε το πρωτόκολλο χαρακτηρισμού δυσκολίας (\S\ref{sec:forensics-protocol})
για να μετρήσουμε την εγγενή διαχωρισιμότητα του συνόλου δεδομένων.

Οι κοινοί γείτονες (\en{CN}) εμφανίζουν πολύ χαμηλή κατανομή: τα θετικά ζεύγη έχουν
μέση τιμή \StructCnMeanPos{} και τα αρνητικά \StructCnMeanNeg{} (εξ ορισμού, εφόσον απουσιάζουν από το γράφημα).
Αντίθετα, η ομοιότητα \en{TF-IDF} παρουσιάζει καθαρή διαχωρισιμότητα: θετικά ζεύγη
έχουν μέση τιμή \StructTfidfMeanPos{} σε σχέση με τα αρνητικά στο \StructTfidfMeanNeg{}.
Χρησιμοποιώντας την
προσέγγιση ποσοστιαίας FPR, επιλέγουμε τα όρια $\theta_{\mathrm{CN}} > \CnThreshold$ και
$\theta_{\mathrm{TF\text{-}IDF}} > \TfidfThreshold$ (1\% FPR στην αρνητική κατανομή).

Ο χαρακτηρισμός αποδίδει τις ακόλουθες κατανομές:

\begin{table}[h!]
  \centering
  \small
  \begin{tabular}{lrcrrc}
    \toprule
    \textbf{Κατηγορία Δυσκολίας} & \textbf{Εκπαίδευση} & & \textbf{Έλεγχος} & \\
    \midrule
    Αυτοβρόχος (τετριμμένο) & \TrainDiffSelfLoopN{} & \TrainDiffSelfLoopPct\% & \TestDiffSelfLoopN{} & \TestDiffSelfLoopPct\% \\
    Υψηλό \en{CN} (τετριμμένο) & \TrainDiffHighCnN{} & \TrainDiffHighCnPct\% & \TestDiffHighCnN{} & 0{,}00\% \\
    Υψηλή ομοιότητα κειμένου & \TrainDiffHighTextsimN{} & \TrainDiffHighTextsimPct\% & \TestDiffHighTextsimN{} & \TestDiffHighTextsimPct\% \\
    Δύσκολο & \TrainDiffHardN{} & \TrainDiffHardPct\% & \TestDiffHardN{} & \TestDiffHardPct\% \\
    \bottomrule
  \end{tabular}
  \caption{Κατανομή δυσκολίας σε σύνολα εκπαίδευσης και ελέγχου.}
  \label{tab:difficulty-breakdown}
\end{table}

\begin{figure}[h!]
  \centering
  \includegraphics[width=0.9\textwidth]{separability_distributions}
  \caption{Κατανομές Κοινών Γειτόνων και ομοιότητας \en{TF-IDF} για θετικά και αρνητικά ζεύγη (χωρίς αυτοβρόχους). Η καθαρή διαχωρισιμότητα της \en{TF-IDF} υποδηλώνει ότι το κειμενικό περιεχόμενο παρέχει ισχυρό σήμα.}
  \label{fig:separability-dist}
\end{figure}

\begin{figure}[h!]
  \centering
  \includegraphics[width=0.7\textwidth]{difficulty_breakdown}
  \caption{Κατανομή της λειτουργικής ετικέτας δυσκολίας στο σύνολο εκπαίδευσης. Η κατηγορία «δύσκολο» δηλώνει ότι δεν ικανοποιείται κανένα από τα τρία απλά κριτήρια.}
  \label{fig:difficulty-breakdown}
\end{figure}

Τα αποτελέσματα δείχνουν ότι το \TrainDiffHardPct\% του συνόλου δεδομένων δεν
ικανοποιεί κανένα από τα τρία απλά κριτήρια, ενώ το \TrainDiffTrivialPct\%
ικανοποιεί τουλάχιστον ένα. Η ονομασία «δύσκολο» είναι λειτουργική και δεν
αποδεικνύει ότι το ζεύγος απαιτεί σημασιολογική συλλογιστική· το τεχνούργημα της
\S\ref{sec:hub-artifact} δείχνει ότι μπορεί να υπάρχει άλλο, μη σχεσιακό σήμα.


\section{Αποτελέσματα \en{CascadeLP} ανά Επίπεδο και Δυσκολία}
\label{sec:cascade-results}

Εκπαιδεύουμε τα τρία επίπεδα ταξινομητών του \en{CascadeLP} αποκλειστικά στο 80\% του
συνόλου εκπαίδευσης (\TrainSplitSize{} ζεύγη, εξαιρουμένων των αυτοβρόχων) και αξιολογούμε σε ένα
παρακρατημένο, στρωματοποιημένο σύνολο επικύρωσης \ValSplitSize{} ζευγών που δεν συμμετέχει
στην εκπαίδευση κανενός ταξινομητή, αλλά έχει επηρεάσει την κατασκευή του
δομικού γραφήματος όπως περιγράφεται στη \S\ref{subsec:data-splits}.
Στα προεπιλεγμένα όρια αξιοπιστίας ($\tau_1 = \TauOneDefault$,
$\tau_2 = \TauTwoDefault$· βλ.\ \S\ref{sec:ablation} για την ευαισθησία ως προς αυτή την επιλογή),
το μοντέλο επιτυγχάνει \en{Macro F1} = \CascadeValFone{} και \en{AUC-ROC} = \CascadeValAuc{} στο σύνολο
επικύρωσης.
Η αυτόνομη βασική γραμμή \en{POS} πετυχαίνει ελαφρώς υψηλότερο \en{Macro F1}
(0{,}9988 έναντι 0{,}9986). Επομένως το αποτέλεσμα του \en{DSAA} δεν δείχνει
βελτίωση του \en{CascadeLP} έναντι κάθε επιμέρους μοντέλου· δείχνει κυρίως ότι
πολλές διαφορετικές μέθοδοι κορεστούν στο ίδιο προβληματικό πεδίο αξιολόγησης.

\begin{table}[h!]
  \centering
  \begin{tabular}{lrrcc}
    \toprule
    \textbf{Επίπεδο} & \textbf{n} & \textbf{Δρομολόγηση} & \textbf{Ακρίβεια} & \textbf{\en{Macro F1}} \\
    \midrule
    Επίπεδο~0 (αυτοβρόχος)      & 0       & 0{,}00\%  & --     & --     \\
    Επίπεδο~1 (\en{Structural}) & \TierOneCount{} & \TierOnePct\% & \TierOneAcc{} & \TierOneFone{} \\
    Επίπεδο~2 (\en{POS + RF})   & \TierTwoCount{} & \TierTwoPct\% & \TierTwoAcc{} & \TierTwoFone{} \\
    Επίπεδο~3 (\en{Embedding})  & \TierThreeCount{} & \TierThreePct\% & \TierThreeAcc{} & \TierThreeFone{} \\
    \bottomrule
  \end{tabular}
  \caption{Ποσοστό δρομολόγησης και απόδοση ανά επίπεδο στο σύνολο επικύρωσης (n=\ValSplitSize{}).}
  \label{tab:cascade-tier-results}
\end{table}

Ο πίνακας~\ref{tab:cascade-tier-results} αποκαλύπτει δύο ευρήματα που δεν προκύπτουν από
τη συνολική βαθμολογία.
Πρώτον, το Επίπεδο~1 έχει ακρίβεια \TierOneAcc{} αλλά \en{Macro F1}
μόλις \TierOneFone --- πρακτικά τυχαία απόδοση παρά την υψηλή ακρίβεια.
Η αιτία είναι η
ακραία ανισορροπία κλάσεων μέσα στο υποσύνολο ζευγών για τα οποία ο
\texttt{\en{StructuralClassifier}} δηλώνει αξιοπιστία $\geq \tau_1$: το \TierOnePosPct\% αυτών
των ζευγών είναι θετικό, και ο ταξινομητής προβλέπει θετική ετικέτα στο 100\% των
περιπτώσεων, χωρίς καμία πρόβλεψη αρνητικής ετικέτας.
Το αποτέλεσμα είναι μηδενική
ανάκληση για την σπάνια αρνητική κλάση, ακριβώς το σενάριο για το οποίο επιλέχθηκε το
\en{Macro F1} ως πρωτεύουσα μετρική (Κεφάλαιο~\ref{ch:experiments}) αντί της απλής
ακρίβειας.
Δεύτερον, το Επίπεδο~3 --- τα \TierThreeCount{} ζεύγη (\TierThreePct\%) που κλιμακώνονται μέχρι τον
\texttt{\en{EmbeddingClassifier}} --- έχει ακρίβεια μόλις \TierThreeAcc.
Η κατεύθυνση του
σφάλματος είναι επίσης ασύμμετρη: οι πραγματικές ετικέτες είναι \TierThreeTruePosPct\% θετικές, ενώ οι
προβλέψεις είναι μόλις \TierThreePredPosPct\% θετικές.
Η ασυμμετρία αυτή αποκαλύπτει συστηματική
μεροληψία: ο ταξινομητής \en{Embedding} τείνει προς υποπρόβλεψη θετικής ετικέτας σε αυτό
το συγκεκριμένο υπόλειμμα ζευγών, που αναλύεται περαιτέρω στη \S\ref{sec:hard-residual}.

Ο πίνακας διασταύρωσης επιπέδου$\times$δυσκολίας (πίνακας~\ref{tab:cascade-tier-difficulty})
δείχνει ότι τα σφάλματα του Επιπέδου~1 και του Επιπέδου~3 συγκεντρώνονται στην κατηγορία
\en{hard}: το Επίπεδο~1 έχει ακρίβεια \TierOneHardAcc{} στα \TierOneHardN{} δύσκολα ζεύγη που χειρίζεται
έναντι \TierOneHighTextsimAcc{} στα τετριμμένα ζεύγη υψηλής ομοιότητας κειμένου που φτάνουν ως αυτό, και το
Επίπεδο~3 καταρρέει στα \TierThreeHardN{} δύσκολα ζεύγη του (ακρίβεια \TierThreeHardAcc) ενώ επιλύει τέλεια τα
υπόλοιπα \TierThreeHighTextsimN{} τετριμμένα ζεύγη που το φτάνουν.
Το Επίπεδο~2 παραμένει σχεδόν αλάνθαστο
(ακρίβεια $\geq$ 0{,}9989) και στις δύο κατηγορίες δυσκολίας, γεγονός που εξηγεί γιατί
χειρίζεται το \TierTwoPct\% του συνόλου επικύρωσης χωρίς να προκαλεί κλιμάκωση.
Από τα \CascadeTotalErrors{}
συνολικά σφάλματα ταξινόμησης στο σύνολο επικύρωσης, το \CascadeErrorsHardPct\% προέρχεται από ζεύγη
που χαρακτηρίζονται \en{hard} και μόλις το \CascadeErrorsTrivialPct\% από τετριμμένα ζεύγη υψηλής ομοιότητας
κειμένου --- τα λάθη του \en{CascadeLP} συγκεντρώνονται σχεδόν αποκλειστικά στο υποσύνολο
που το πρωτόκολλο χαρακτηρισμού δυσκολίας (\S\ref{sec:forensics-protocol}) είχε ήδη
προσδιορίσει ως μη διαχωρίσιμο από τα απλά κριτήρια.

\begin{table}[h!]
  \centering
  \small
  \begin{tabular}{lrrr}
    \toprule
    \textbf{Επίπεδο} & \textbf{\en{hard}} & \textbf{Υψηλό \en{CN}} & \textbf{Υψηλή Ομοιότητα Κειμένου} \\
    \midrule
    Επίπεδο~1 & n=\TierOneHardN, \TierOneHardAcc{} & n=\TierOneHighCnN, 1{,}0 & n=\TierOneHighTextsimN, 1{,}0 \\
    Επίπεδο~2 & n=\TierTwoHardN, \TierTwoHardAcc{} & ---         & n=\TierTwoHighTextsimN, \TierTwoHighTextsimAcc{} \\
    Επίπεδο~3 & n=\TierThreeHardN, \TierThreeHardAcc{} & ---         & n=\TierThreeHighTextsimN, 1{,}0 \\
    \bottomrule
  \end{tabular}
  \caption{Πίνακας διασταύρωσης επιπέδου$\times$δυσκολίας: πλήθος ζευγών και ακρίβεια ανά
  συνδυασμό, σύνολο επικύρωσης.}
  \label{tab:cascade-tier-difficulty}
\end{table}

Στο σύνολο ελέγχου, όπου δεν υπάρχουν πραγματικές ετικέτες στα δεδομένα μας και άρα
μετράμε τοπικά μόνο ποσοστά κλήσεων ανά επίπεδο, παρατηρούμε μια αξιοσημείωτη μετατόπιση
σε σχέση με το σύνολο επικύρωσης: το Επίπεδο~1 χειρίζεται μόλις το \TestTierOnePct\% των ζευγών
ελέγχου (\TestTierOneCount{} από \TestPairs{}) έναντι \TierOnePct\% στο σύνολο επικύρωσης, με το Επίπεδο~2 να
απορροφά την υπόλοιπη κίνηση στο \TestTierTwoPct\%.
Το Επίπεδο~0 χειρίζεται \SelfLoopsTest{} ζεύγη (\TestTierZeroPct\%)
--- αριθμός που ταυτίζεται ακριβώς με το πλήθος αυτοβρόχων που καταγράφηκε στον έλεγχο
\en{data leakage} (\S\ref{sec:leakage-results}), λειτουργώντας ως ανεξάρτητος έλεγχος
ορθότητας της διαδρομής δρομολόγησης.
Η μετατόπιση στο Επίπεδο~1 δεν είναι σφάλμα
υλοποίησης, αλλά άμεση συνέπεια του τρόπου κατασκευής του δομικού γραφήματος: το γράφημα
περιλαμβάνει αποκλειστικά τις θετικές ακμές του συνόλου \textit{εκπαίδευσης}
(\S\ref{sec:cascade-routing}), οπότε και οι δύο κόμβοι ενός ζεύγους εμφανίζονται σε αυτό
για το \GraphCoverageTrainPct\% των ζευγών εκπαίδευσης/επικύρωσης, αλλά μόλις για το \GraphCoverageTestPct\% των ζευγών
ελέγχου.
Τα ζεύγη ελέγχου είναι δηλαδή δομικά πιο «ψυχρά» από τα ζεύγη επικύρωσης, το
Επίπεδο~1 στερείται σήματος γι' αυτά και η δρομολόγηση κλιμακώνει σχεδόν αυτόματα στο
Επίπεδο~2.
Η ανάλυση της απόδοσης ειδικά στο υποσύνολο μηδενικού \en{CN}
αναλαμβάνεται στη \S\ref{sec:cold-start-results}.

Παρότι το \en{DSAA 2023} δεν δημοσιεύει τις πραγματικές ετικέτες του συνόλου ελέγχου,
ο διαγωνισμός διατηρεί ενεργό πίνακα κατάταξης στο \en{Kaggle} που βαθμολογεί
υποβληθείσες προβλέψεις.
Υποβάλλοντας τις προβλέψεις του \en{CascadeLP} στο σύνολο
ελέγχου (\texttt{\en{outputs\slash{}predictions\slash{}cascade\_submission.csv}}) λάβαμε δημόσια
βαθμολογία \en{Macro F1} = \KaggleBaselinePublic{}. Η συμφωνία με τη βαθμολογία
επικύρωσης (\CascadeValFone{}) δείχνει ότι το ίδιο μοτίβο είναι παρόν και στο
κρυφό σύνολο ελέγχου, αλλά δεν διαχωρίζει τη σχεσιακή πρόβλεψη από την αξιοποίηση
του τεχνουργήματος δειγματοληψίας της \S\ref{sec:hub-artifact}.


\section{Αποτελέσματα στο Ανεξάρτητο Σύνολο \en{Wikipedia}}
\label{sec:wiki-fresh}
\label{subsec:wiki-fresh-results}

Στο δεύτερο σύνολο επανεκπαιδεύουμε από την αρχή την ίδια διαμόρφωση του
\en{CascadeLP}, χωρίς αλλαγή αρχιτεκτονικής ή ορίων, και αξιολογούμε τους ίδιους
πέντε αυτόνομους ταξινομητές. Τα αποτελέσματα παρουσιάζονται στον
Πίνακα~\ref{tab:wiki-fresh-results}.

\begin{table}[h!]
  \centering
  \begin{tabular}{lrr}
    \toprule
    \textbf{Μοντέλο} & \textbf{\en{Macro F1}} & \textbf{\en{AUC-ROC}} \\
    \midrule
    Δομικό (\en{CN/Jaccard/AA/PA})       & \WfStructuralFone{} & \WfStructuralAuc{} \\
    \en{TF-IDF + Λογιστική Παλινδρόμηση} & \WfTfidfFone{}      & \WfTfidfAuc{} \\
    \en{POS + Random Forest}             & \WfPosFone{}        & \WfPosAuc{} \\
    Ομοιότητα \en{Sentence-Transformer}  & \WfEmbeddingFone{}  & \WfEmbeddingAuc{} \\
    \en{TF-IDF + SVM}                    & \WfSvmFone{}        & \WfSvmAuc{} \\
    \en{CascadeLP}                       & \textbf{\WfCascadeFone{}} & \textbf{\WfCascadeAuc{}} \\
    \bottomrule
  \end{tabular}
  \caption{Αποτελέσματα στο ανεξάρτητο σύνολο Απλής Αγγλικής \en{Wikipedia}
  (n=\WfValSplitSize{} ζεύγη επικύρωσης). Όλα τα μοντέλα επανεκπαιδεύονται από
  την αρχή χωρίς αλλαγή της αρχιτεκτονικής ή των υπερπαραμέτρων τους.}
  \label{tab:wiki-fresh-results}
\end{table}

Το \en{CascadeLP} πετυχαίνει \en{Macro F1} = \WfCascadeFone{}, υψηλότερο από
κάθε μεμονωμένη οικογένεια χαρακτηριστικών. Ο δομικός ταξινομητής είναι η
ισχυρότερη αυτόνομη οικογένεια με \WfStructuralFone{}, γεγονός που επιβεβαιώνει
ότι το πυκνότερο γράφημα παρέχει χρήσιμο σχεσιακό σήμα. Η δρομολόγηση ακολουθεί
επίσης τον σχεδιαστικό στόχο: το Επίπεδο~1 επιλύει \WfTierOneCount{} ζεύγη
(\WfTierOnePct\%), το Επίπεδο~2 \WfTierTwoCount{} (\WfTierTwoPct\%) και το
Επίπεδο~3 \WfTierThreeCount{} (\WfTierThreePct\%). Τα ποσοστά αυτά μετρούν ποιο
επίπεδο έλαβε την τελική απόφαση πάνω σε προϋπολογισμένα χαρακτηριστικά· δεν
αποτελούν μέτρηση ίσης ποσοστιαίας μείωσης του κόστους κωδικοποίησης.


\section{Σύγκριση με Μοντέλο Αναφοράς \en{SVM}}
\label{sec:svm-results}

Ως μοντέλο αναφοράς εκπαιδεύουμε έναν ταξινομητή \en{TF-IDF + SVM} (RBF kernel SVC
σε ένα στρωματοποιημένο υποσύνολο \SvmSampleSize{} ζευγών για να αποφύγουμε το τετραγωνικό υπολογιστικό κόστος
σε 948K ζεύγη). Σημειώνεται ότι ο περιορισμός του μεγέθους εκπαίδευσης είναι εσκεμμένη
επιλογή λόγω του απαγορευτικού κόστους του αλγορίθμου. Συνεπώς, η σύγκριση με τα 
άλλα μοντέλα γίνεται με άξονα τον λόγο κόστους-απόδοσης σε ένα πρακτικά
εφικτό σενάριο εκπαίδευσης, και δεν αντικατοπτρίζει απαραίτητα τη μέγιστη
θεωρητική ικανότητα της μεθόδου \en{SVM}.
Τα αποτελέσματα επικύρωσης είναι:

\begin{table}[h!]
  \centering
  \begin{tabular}{lcc}
    \toprule
    \textbf{Μέτρο} & \textbf{Τιμή} \\
    \midrule
    \en{Macro F1-score} & \SvmFone{} \\
    \en{AUC-ROC} & \SvmAuc{} \\
    \en{Macro F1} μηδενικού \en{CN} & \SvmCsFone{} \\
    Καθυστέρηση & \SvmLatency{} ms/ζεύγος \\
    \bottomrule
  \end{tabular}
  \caption{\en{SVM} αποτελέσματα επικύρωσης σε \ValSplitSize{} ζεύγη.}
  \label{tab:svm-results}
\end{table}

\begin{figure}[h!]
  \centering
  \includegraphics[width=0.8\textwidth]{svm_metrics}
  \caption{Απόδοση μοντέλου αναφοράς \en{TF-IDF + SVM} σε σύνολο επικύρωσης. Το \en{SVM} επιτυγχάνει \en{Macro F1} = \SvmFone{} και \en{AUC-ROC} = \SvmAuc{}, υποδηλώνοντας την ισχυρή προβλεπτική ικανότητα της κειμενικής ομοιότητας.}
  \label{fig:svm-metrics}
\end{figure}

Το \en{SVM} επιτυγχάνει υψηλή απόδοση σε όλες τις μετρικές, υπογραμμίζοντας ότι η
\en{TF-IDF} ομοιότητα του κειμένου είναι ένας ισχυρός μεμονωμένος δείκτης για τη
πρόβλεψη ακμών σε αυτό το σύνολο δεδομένων.
Η σύγκριση με τη \en{CascadeLP}
θα δείξει κατά πόσον οι πολυστοιχειακές δομικές και σημασιολογικές προσεγγίσεις
παρέχουν περαιτέρω κέρδη απόδοσης ή απλώς αναδημιουργούν τα αποτελέσματα που
καταδεικνύονται ήδη από τη \en{TF-IDF}.
